% Luc Destombe --- licence CC BY-NC-SA 4.0
% https://creativecommons.org/licenses/by-nc-sa/4.0/deed.fr
% Fichier autonome : compiler avec pdflatex, deux fois (signets, nombre de pages).
\documentclass[11pt,a4paper]{article}
\makeatletter
\newif\iflycee@unecolonne
\lycee@unecolonnetrue
\newif\iflycee@palatino
\lycee@palatinotrue

\RequirePackage[utf8]{inputenc}
\RequirePackage[T1]{fontenc}
\RequirePackage[french]{babel}
\RequirePackage{etoolbox}

\newcommand{\ShorthandsOff}{\@ifundefined{shorthandoff}{}{\shorthandoff{;:!?}}}
\newcommand{\ShorthandsOn}{\@ifundefined{shorthandon}{}{\shorthandon{;:!?}}}
\ShorthandsOff
\AtBeginDocument{\ShorthandsOn}
\AtBeginEnvironment{tikzpicture}{\ShorthandsOff}

\RequirePackage{geometry}
\RequirePackage{fancyhdr}
\RequirePackage{lastpage}
\RequirePackage{multicol}
\RequirePackage{array}
\RequirePackage{enumitem}
\RequirePackage{xcolor}
\RequirePackage{needspace}

\RequirePackage{amsmath,amssymb}
\RequirePackage{mathpazo}
\DeclareMathSizes{10.95}{10.95}{8}{5.5}
\begingroup\catcode`\_=\active \catcode`\^=\active
\gdef_#1{\sb{\scriptscriptstyle #1}}
\gdef^#1{\sp{\scriptscriptstyle\lycee@moinsepais #1}}
\endgroup
\newcommand\lycee@moins{\mathbin{%
  \setbox\z@\hbox{$\m@th\scriptscriptstyle\lycee@moinsorig$}%
  \dimen@=.55\wd\z@ \advance\dimen@-.5pt
  \hbox to.75\wd\z@{\hss$\m@th\scriptscriptstyle\vcenter{\hbox{%
    \tikz\draw[line width=.5pt,line cap=round](0,0)--(\the\dimen@,0);}}$\hss}}}
\begingroup\lccode`\~=`\- \lowercase{\endgroup\let~}\lycee@moins
\newcommand\lycee@moinsepais{\mathcode`\-="8000 }
\AtBeginDocument{\mathchardef\lycee@moinsorig=\mathcode`\-\relax}
\AtBeginDocument{\catcode`\_=12 \mathcode`\_="8000
  \catcode`\^=12 \mathcode`\^="8000 }
\newcommand\lycee@exposants{%
  \fontdimen13\textfont2=.48\fontdimen6\textfont2
  \fontdimen14\textfont2=.48\fontdimen6\textfont2
  \fontdimen15\textfont2=.48\fontdimen6\textfont2
  \fontdimen13\scriptfont2=.48\fontdimen6\scriptfont2
  \fontdimen14\scriptfont2=.48\fontdimen6\scriptfont2
  \fontdimen15\scriptfont2=.48\fontdimen6\scriptfont2}
\every@math@size\expandafter{\the\every@math@size\lycee@exposants}
\newlength{\retraitcalculs}
\setlength{\retraitcalculs}{2em}

\RequirePackage{tikz}
\usetikzlibrary{arrows.meta,calc,babel,shapes.misc}
\RequirePackage[most]{tcolorbox}
\tcbuselibrary{skins,breakable}

\definecolor{coulDef}{HTML}{1F4E79}
\definecolor{coulProp}{HTML}{7030A0}
\definecolor{coulRem}{HTML}{C55A11}
\definecolor{coulEx}{HTML}{2E7D32}
\definecolor{coulExo}{HTML}{B03A2E}
\definecolor{coulPart}{HTML}{1F4E79}
\definecolor{coulMeth}{HTML}{1B6E4F}
\definecolor{coulCourbe}{HTML}{0B5ED7}
\definecolor{coulCourbeB}{HTML}{C00000}
\definecolor{coulCourbeC}{HTML}{2E7D32}
\definecolor{coulGrille}{HTML}{8C8C8C}
\definecolor{coulRep}{HTML}{1B6E4F}
\definecolor{coulBar}{HTML}{2E7D32}

\newcommand{\R}{\mathbb{R}}
\newcommand{\Rs}{\mathbb{R}^{*}}
\renewcommand{\le}{\leqslant}
\renewcommand{\ge}{\geqslant}
\renewcommand{\approx}{\simeq}
\newcommand{\vect}[1]{\overrightarrow{#1}}
\newcommand{\norme}[1]{\left\lVert #1 \right\rVert}
\newcommand{\intervalle}[2]{\left[#1\,;#2\right]}
\RequirePackage{cancel}
\renewcommand{\CancelColor}{\color{coulExo}}
\newcommand{\fois}[1]{\times{\color{coulExo}#1}}
\newcommand{\barre}[1]{\cancel{#1}}

\tikzset{grille/.style={coulGrille,line width=0.4pt}}
\newcommand{\quadrillage}[4]{\draw[grille,step=1] (#1,#2) grid (#3,#4);}
\tikzset{
  croix/.style={cross out,draw,line width=0.6pt,inner sep=0pt,minimum size=4pt},
  croixdroite/.style={croix,rotate=45,line width=0.8pt},
}
\newcommand{\pt}[3]{\path (#1) node[croix]{} node[#2,font=\small,inner sep=2.5pt]{$#3$};}

\newcommand{\lycee@repere}[4]{%
  \draw[grille,step=1] (#1,#2) grid (#3,#4);
  \draw[-{Stealth[length=2mm]},thick] (#1,0)--({#3+0.4},0) node[below left=-1pt]{$x$};
  \draw[-{Stealth[length=2mm]},thick] (0,#2)--(0,{#4+0.4}) node[below left=-1pt]{$y$};
  \node[below left,font=\scriptsize,inner sep=1.5pt] at (0,0) {$0$};}
\newcommand{\lycee@gradx}[1]{\foreach \k/\l in {#1}{%
  \draw (\k,0.07)--(\k,-0.07) node[below,font=\scriptsize,inner sep=1.5pt]{$\l$};}}
\newcommand{\lycee@grady}[1]{\foreach \k/\l in {#1}{%
  \draw (0.07,\k)--(-0.07,\k) node[left,font=\scriptsize,inner sep=1.5pt]{$\l$};}}
\newcommand{\lycee@intff}[2]{\left[#1\,;#2\right]}
\newcommand{\lycee@intfo}[2]{\left[#1\,;#2\right[}
\newcommand{\lycee@intof}[2]{\left]#1\,;#2\right]}
\newcommand{\lycee@intoo}[2]{\left]#1\,;#2\right[}
\newcommand{\lycee@ens}[1]{\left\{#1\right\}}
\AtBeginDocument{%
  \providecommand{\repere}{\lycee@repere}%
  \providecommand{\gradx}{\lycee@gradx}%
  \providecommand{\grady}{\lycee@grady}%
  \providecommand{\Cf}{\mathcal{C}\sb{\scriptscriptstyle f}}%
  \providecommand{\Cg}{\mathcal{C}\sb{\scriptscriptstyle g}}%
  \providecommand{\intff}{\lycee@intff}%
  \providecommand{\intfo}{\lycee@intfo}%
  \providecommand{\intof}{\lycee@intof}%
  \providecommand{\intoo}{\lycee@intoo}%
  \providecommand{\ens}{\lycee@ens}}

\newlist{questions}{enumerate}{3}
\setlist[questions,1]{label=\arabic*\textdegree),leftmargin=*,itemsep=2pt,topsep=3pt}
\setlist[questions,2]{label=\alph*),leftmargin=*,itemsep=1pt,topsep=2pt}
\setlist[questions,3]{label=\roman*),leftmargin=*,itemsep=1pt,topsep=1pt}
\setlist[itemize]{label=\textbullet,leftmargin=*,itemsep=2pt,topsep=3pt}

\newcommand{\besoinplace}[1]{\par\penalty\@M\begingroup
  \dimen@=#1\relax\dimen@ii=\pagegoal\advance\dimen@ii-\pagetotal
  \ifdim\dimen@>\dimen@ii\ifdim\dimen@ii>\z@\vfil\fi\break\fi\endgroup}

\newcommand{\lycee@pied}{\thepage/\pageref{LastPage}}
\newcommand{\entete}[2]{%
  \pagestyle{fancy}\fancyhf{}%
  \fancyhead[L]{\footnotesize\color{coulPart}\textbf{#1}}%
  \fancyhead[R]{\footnotesize\color{coulPart}\textbf{#2}}%
  \fancyfoot[C]{\footnotesize\lycee@pied}%
  \renewcommand{\headrulewidth}{0.4pt}%
  \renewcommand{\headrule}{\hbox to\headwidth{%
    \color{coulPart!40}\leaders\hrule height \headrulewidth\hfill}}}

\RequirePackage{ccicons}
\newcommand{\lycee@licence}{{\scriptsize\color{gray}%
  \copyright\ Luc Destombe --- \ccbyncsa\ CC BY-NC-SA 4.0}}
\newif\iflycee@licence \lycee@licencetrue
\newcommand{\sanslicence}{\lycee@licencefalse}
\AtBeginDocument{\iflycee@licence\AddToHookNext{shipout/foreground}{%
  \put(0,0){\rlap{\kern\dimexpr1in+\hoffset+\oddsidemargin+\textwidth\relax
    \llap{\raisebox{-\dimexpr1in+\voffset+\topmargin+\headheight+\headsep
      +\textheight+\footskip\relax}{\lycee@licence}}}}}\fi}

\newcommand{\formulesserrees}{%
  \setlength{\abovedisplayskip}{3pt}\setlength{\belowdisplayskip}{3pt}%
  \setlength{\abovedisplayshortskip}{2pt}\setlength{\belowdisplayshortskip}{3pt}}

\setlength{\parindent}{0pt}

\RequirePackage[implicit=false,hidelinks,pdfencoding=auto,psdextra,bookmarksopen,
  bookmarksopenlevel=1,pdfcreator={},pdfproducer={}]{hyperref}
\RequirePackage{bookmark}
\hypersetup{pdfauthor={Luc Destombe},pdfkeywords={CC BY-NC-SA 4.0}}
\newcounter{lycee@signet}
\newcommand{\signet}[3][]{\stepcounter{lycee@signet}%
  \edef\lycee@dest{\if\relax\detokenize{#1}\relax
    signet.\arabic{lycee@signet}\else#1\fi}%
  \hypertarget{\lycee@dest}{}\bookmark[level=#2,dest=\lycee@dest]{#3}}
\pdfstringdefDisableCommands{%
  \def\R{R}\def\vect#1{#1}\def\Cf{Cf}\def\Cg{Cg}\def\fois#1{×#1}%
  \def\barre#1{#1}\def\intervalle#1#2{[#1;#2]}\def\intff#1#2{[#1;#2]}%
  \def\textsuperscript#1{#1}\def\dfrac#1#2{#1/#2}\def\frac#1#2{#1/#2}%
  \def\vec#1{#1⃗}}%
\RequirePackage[official]{eurosym}

\tcbset{exercice corrige/.style={
  enhanced,breakable,before upper=\raggedright,
  colback=white,colframe=coulExo,
  boxrule=0.8pt,arc=2pt,
  left=6pt,right=6pt,top=8pt,bottom=5pt,
  before skip=9pt,after skip=4pt,
  coltitle=white,fonttitle=\bfseries\small,
  attach boxed title to top left={xshift=8pt,yshift=-\tcboxedtitleheight/2},
  boxed title style={colback=coulExo,arc=2pt,boxrule=0pt}}}
\newcounter{exo}
\newtcolorbox{exercice}[1][]{exercice corrige,
  phantom={\signet[exercice-\arabic{exo}]{2}{Exercice \arabic{exo}}},
  title={Exercice \theexo\ifx\relax#1\relax\else{} --- #1\fi}}
\newenvironment{exo}[1][\relax]{\refstepcounter{exo}\begin{exercice}[#1]}{\end{exercice}}

\geometry{top=1.8cm,bottom=1cm,left=1cm,right=1cm,headheight=14pt,footskip=0.6cm}
\renewcommand{\lycee@pied}{\thepage}

\renewcommand{\theexo}{\ifnum\value{exo}<10 0\fi\arabic{exo}}
\tcbset{exercice corrige/.style={
  enhanced,breakable,before upper=\raggedright,
  colback=white,colframe=coulExo,
  boxrule=0.9pt,arc=2pt,
  left=8pt,right=8pt,top=8pt,bottom=6pt,
  before skip=8pt,after skip=6pt,
  coltitle=white,fonttitle=\bfseries,
  attach boxed title to top left={xshift=8pt,yshift=-\tcboxedtitleheight/2},
  boxed title style={colback=coulExo,arc=2pt,boxrule=0pt}}}

\newcommand{\entetefiche}[2]{%
  \begin{center}
    \begin{tcolorbox}[enhanced,width=0.92\linewidth,colback=white,colframe=coulPart,
        boxrule=1.4pt,arc=3pt,halign=center,top=6pt,bottom=6pt]
      {\large\bfseries\color{coulPart} CORRIGÉ --- FICHE D'EXERCICES}\\[3pt]
      {\normalsize\bfseries #1}\\[2pt]
      {\normalsize #2}
    \end{tcolorbox}
  \end{center}}

\newcounter{partie}
\newcommand{\partie}[1]{%
  \stepcounter{partie}%
  \besoinplace{8\baselineskip}\vspace{4pt}%
  \begin{tcolorbox}[enhanced,colback=coulPart!8,colframe=coulPart,boxrule=0pt,
      leftrule=3pt,arc=0pt,left=6pt,right=4pt,top=3pt,bottom=3pt,
      before skip=6pt,after skip=2pt]
    \bfseries\color{coulPart}\Roman{partie} --- #1\signet{1}{\Roman{partie} --- #1}
  \end{tcolorbox}}
\newcommand{\partiesansnum}[1]{%
  \besoinplace{8\baselineskip}\vspace{4pt}%
  \begin{tcolorbox}[enhanced,colback=coulPart!8,colframe=coulPart,boxrule=0pt,
      leftrule=3pt,arc=0pt,left=6pt,right=4pt,top=3pt,bottom=3pt,
      before skip=6pt,after skip=2pt]
    \bfseries\color{coulPart}#1\signet{1}{#1}
  \end{tcolorbox}}

\setlist[questions,1]{label=\arabic*\textdegree),leftmargin=*,itemsep=3pt,topsep=3pt}
\setlist[questions,2]{label=\alph*),leftmargin=*,itemsep=2pt,topsep=2pt}
\setlist[questions,3]{label=\roman*),leftmargin=*,itemsep=1pt,topsep=1pt}
\newlist{expr}{enumerate}{1}
\setlist[expr]{label={\Alph*\ $=$},leftmargin=*,itemsep=2pt,topsep=3pt}

\newcommand{\res}[1]{{\color{coulRep}#1}}
\newcommand{\rem}[1]{\par\smallskip{\small\itshape\color{black!60}$\rightarrow$ #1}}
\newenvironment{calculs}{\par\smallskip\noindent$\displaystyle\begin{aligned}[t]}%
  {\end{aligned}$\par\smallskip}
\newcommand{\justif}[1]{\quad\text{\small\itshape\color{black!60}#1}}
\newcommand{\cas}[1]{\par\smallskip\noindent\textbf{\color{coulExo}#1}\nobreak}
\newcommand{\explique}[2]{\par\smallskip\noindent
  \begin{minipage}[t]{0.6\linewidth}\raggedright #1\end{minipage}\hfill
  \begin{minipage}[t]{0.37\linewidth}\raggedright\leavevmode
    {\small\itshape\color{black!60}#2}\end{minipage}\par}

\setlength{\parskip}{2pt}

\makeatother
\usepackage{tkz-tab}
\entete{Chapitre 2 --- Fonctions affines et pourcentages --- corrigé}{1\textsuperscript{re} mathématiques spécifiques}
\setlength{\parindent}{0pt}
\setlength{\jot}{5pt}

\let\casligne\cas
\renewcommand{\cas}[1]{\casligne{#1}\par\nopagebreak}
\newcommand{\etiq}[1]{\textbf{\color{coulExo}#1}\ }

\newcommand{\cotefig}[3][0.36]{\par\noindent
  \begin{minipage}[t]{\dimexpr0.97\linewidth-#1\linewidth\relax}\raggedright #2\end{minipage}\hfill
  \begin{minipage}[t]{#1\linewidth}\vspace{0pt}\centering #3\end{minipage}\par}

\tikzset{
  lecture/.style={densely dashed,coulRep,line width=0.8pt},
  pas/.style={densely dashed,line width=0.9pt},
}
\newcommand{\reperefiche}{%
  \quadrillage{-4}{-4}{5}{5}
  \draw[->,>=Stealth,black!70] (-4,0) -- (5.5,0) node[below,font=\scriptsize] {$x$};
  \draw[->,>=Stealth,black!70] (0,-4) -- (0,5.5) node[left,font=\scriptsize] {$y$};
  \foreach \i in {-4,...,-1,1,2,...,4} {\node[below,font=\tiny,inner sep=1.5pt] at (\i,0) {$\i$};}
  \foreach \i in {-4,...,-1,1,2,...,4} {\node[left,font=\tiny,inner sep=1.5pt] at (0,\i) {$\i$};}
  \node[below left,font=\tiny,inner sep=1.5pt] at (0,0) {$0$};}
\newcommand{\signes}[4]{%
  \begin{tikzpicture}
  \tkzTabInit[lgt=1.4,espcl=1.5]{$x$/0.7,$#1$/0.7}{$-\infty$,$#2$,$+\infty$}
  \tkzTabLine{,#3,z,#4,}
  \end{tikzpicture}}

\begin{document}

\entetefiche{Chapitre 2 --- Fonctions affines et pourcentages}{1\textsuperscript{re} mathématiques spécifiques}

\partie{Appliquer un pourcentage}

\begin{exo}[Calcul mental]
\explique{On passe par un pourcentage facile ($10\,\%$, $1\,\%$, $25\,\%$, $50\,\%$),
puis on multiplie.}{tableau « Pour calculer de tête » du cours}
\explique{\etiq{a)} $10\,\%$ de $420$ : $420\div 10=\res{42}$.}{$10\,\%$ : on divise
par $10$}
\explique{\etiq{b)} $25\,\%$ de $60$ : le quart de $60$, soit $\res{15}$.}{$25\,\%
=\dfrac{1}{4}$}
\explique{\etiq{c)} $50\,\%$ de $48$ : la moitié de $48$, soit $\res{24}$.}{$50\,\%
=\dfrac{1}{2}$}
\explique{\etiq{d)} $1\,\%$ de $900$ : $900\div 100=\res{9}$.}{$1\,\%$ : on divise par
$100$}
\explique{\etiq{e)} $10\,\%$ de $80$ font $8$, donc $5\,\%$ en font $\res{4}$.}{$5\,\%$ :
la moitié de $10\,\%$}
\explique{\etiq{f)} $25\,\%$ de $20$ font $5$, donc $75\,\%$ en font $3\times 5=\res{15}$.}{$75\,\%$,
c'est $3$ fois $25\,\%$}
\explique{\etiq{g)} $10\,\%$ de $35$ font $3{,}5$, donc $20\,\%$ en font
$2\times 3{,}5=\res{7}$.}{$20\,\%$, c'est $2$ fois $10\,\%$}
\explique{\etiq{h)} $10\,\%$ de $250$ font $25$, donc $30\,\%$ en font
$3\times 25=\res{75}$.}{$30\,\%$, c'est $3$ fois $10\,\%$}
\explique{\etiq{i)} $10\,\%$ de $15$ font $1{,}5$, donc $40\,\%$ en font
$4\times 1{,}5=\res{6}$.}{$40\,\%$, c'est $4$ fois $10\,\%$}
\rem{Contrôle de bon sens : $40\,\%$ est moins que la moitié ; le résultat de i) doit
donc être inférieur à la moitié de $15$, c'est-à-dire à $7{,}5$. C'est le cas.}
\end{exo}

\begin{exo}[Trois écritures d'un même nombre]
\begin{center}
\renewcommand{\arraystretch}{2.4}
\begin{tabular}{|l|c|c|c|c|c|}
\hline
Fraction & $\dfrac{3}{5}$ & $\res{\dfrac{9}{100}}$ & $\res{\dfrac{9}{20}}$ & $\dfrac{1}{4}$ & $\res{\dfrac{3}{2}}$\\
\hline
Décimal & $\res{0{,}6}$ & $0{,}09$ & $\res{0{,}45}$ & $\res{0{,}25}$ & $1{,}5$\\
\hline
Pourcentage & $\res{60\,\%}$ & $\res{9\,\%}$ & $45\,\%$ & $\res{25\,\%}$ & $\res{150\,\%}$\\
\hline
\end{tabular}
\end{center}
\explique{L'outil est la définition du cours : $t\,\%=\dfrac{t}{100}$. On cherche donc
chaque fois un nombre de centièmes.}{« pour cent » : sur $100$}
\cas{1\textsuperscript{re} colonne}
\begin{calculs}
\dfrac{3}{5} &= \dfrac{3\fois{20}}{5\fois{20}} &&\justif{on veut $100$ en bas : $5\times 20=100$}\\
\dfrac{3}{5} &= \dfrac{60}{100} &&\justif{$60$ centièmes}
\end{calculs}
\explique{Soit $0{,}6$ en écriture décimale, et $60\,\%$.}{}
\cas{2\textsuperscript{e} colonne}
\explique{$0{,}09$, c'est $9$ centièmes : $\dfrac{9}{100}$, soit $9\,\%$.}{le $9$ est au
rang des centièmes}
\cas{3\textsuperscript{e} colonne}
\begin{calculs}
45\,\% &= \dfrac{45}{100} &&\justif{définition du pourcentage}\\
45\,\% &= \dfrac{9\times\barre{5}}{20\times\barre{5}} &&\justif{$45=9\times 5$ et
$100=20\times 5$}\\
45\,\% &= \dfrac{9}{20} &&
\end{calculs}
\explique{En écriture décimale : $\dfrac{45}{100}=0{,}45$.}{}
\cas{4\textsuperscript{e} colonne}
\explique{$\dfrac{1}{4}=\dfrac{25}{100}$ : soit $0{,}25$, et $25\,\%$.}{on multiplie en haut
et en bas par $25$}
\cas{5\textsuperscript{e} colonne}
\explique{$1{,}5=\dfrac{150}{100}$ : soit $150\,\%$.}{$1{,}5$, c'est $150$ centièmes}
\explique{Et $1{,}5=\dfrac{15}{10}=\dfrac{3}{2}$.}{on divise en haut et en bas par $5$}
\rem{Erreur fréquente : écrire $0{,}09=90\,\%$. Or $90\,\%=0{,}9$ ; $0{,}09$ ne fait que
$9$ centièmes. Un pourcentage peut dépasser $100\,\%$ : $1{,}5=150\,\%$.}
\end{exo}

\begin{exo}[Retrouver le total]
\explique{Ici, on connaît une partie et son pourcentage ; on cherche le total, qui
vaut $100\,\%$.}{le total est plus grand que la partie}
\cas{a)}
\explique{$5\,\%$ est la vingtième partie de $100\,\%$.}{$20\times 5=100$}
\explique{Le total est donc $20$ fois plus grand que la partie : $20\times 30=600$.}{}
\explique{La ferme compte $\res{600}$ animaux.}{contrôle : $10\,\%$ de $600$ font $60$,
et $5\,\%$ en font $30$}
\cas{b)}
\explique{$30\,\%$ font $9$ élèves, donc $10\,\%$ font $9\div 3=3$ élèves.}{$10\,\%$, c'est le
tiers de $30\,\%$}
\explique{$100\,\%$ font $10\times 3=30$ : la classe compte $\res{30}$ élèves.}{$100\,\%$, c'est
$10$ fois $10\,\%$}
\cas{c)}
\explique{$10\,\%$ du prix font $8$~€, donc le prix entier vaut $10\times 8=80$ : le jeu
coûte $\res{80}$~€.}{$100\,\%$, c'est $10$ fois $10\,\%$}
\cas{d)}
\explique{$75\,\%$ font $60$ adhérents, donc $25\,\%$ font $60\div 3=20$.}{$75\,\%$,
c'est $3$ fois $25\,\%$}
\explique{$100\,\%$ font $4\times 20=80$ : le club compte $\res{80}$ adhérents.}{$100\,\%$,
c'est $4$ fois $25\,\%$}
\rem{Erreur fréquente : calculer $5\,\%$ de $30$ en a). On ne cherche pas une partie de
$30$ : $30$ est déjà la partie, et le total est plus grand.}
\end{exo}

\begin{exo}[Augmenter ou diminuer]
\explique{Comme dans la méthode 2 : on calcule d'abord le montant de la hausse ou de la
baisse ; une hausse s'ajoute, une baisse se retire.}{}
\cas{a)}
\explique{$10\,\%$ de $60$ font $6$ ; $5\,\%$ en font $3$ ; donc $15\,\%$ en font
$6+3=9$.}{$15\,\%=10\,\%+5\,\%$ : le montant de la hausse}
\explique{Nouveau prix : $60+9=\res{69}$~€.}{une hausse s'ajoute}
\cas{b)}
\explique{$10\,\%$ de $900$ font $90$ ; $20\,\%$ en font $180$.}{« soldé à $-20\,\%$ » :
baisse de $20\,\%$}
\explique{Prix soldé : $900-180=\res{720}$~€.}{une baisse se retire}
\cas{c)}
\explique{$1\,\%$ de $2\,000$ font $20$ ; $3\,\%$ en font $60$.}{$1\,\%$ : on divise par
$100$}
\explique{Nouveau salaire : $2\,000+60=\res{2\,060}$~€.}{}
\cas{d)}
\explique{$10\,\%$ de $6\,000$ font $600$ ; $5\,\%$ en font $300$.}{la moitié de
$10\,\%$}
\explique{La commune compte ensuite $6\,000-300=\res{5\,700}$ habitants.}{une baisse se
retire}
\rem{Erreur fréquente : répondre $9$~€ en a). $9$~€ est le montant de la hausse, pas le
nouveau prix.}
\end{exo}

\begin{exo}[Prix hors taxe, prix TTC]
\cas{a)}
\explique{$10\,\%$ de $350$ font $35$ ; la TVA, $20\,\%$, vaut $2\times 35=\res{70}$~€.}{la
TVA est un pourcentage du prix hors taxe}
\explique{Prix TTC : $350+70=\res{420}$~€.}{la taxe s'ajoute au prix hors taxe}
\cas{b)}
\explique{La baisse porte sur le prix TTC, $420$~€ : $10\,\%$ de $420$ font $42$.}{c'est
le prix affiché en magasin qui est soldé}
\explique{Prix soldé : $420-42=\res{378}$~€.}{}
\rem{Erreur fréquente : calculer $10\,\%$ de $350$. Chaque pourcentage porte sur la
valeur du moment : d'abord $350$~€, puis $420$~€.}
\end{exo}

\partie{Calculer un pourcentage}

\begin{exo}[Proportions]
\explique{La proportion est $p=\dfrac{\text{partie}}{\text{total}}$ ; pour l'écrire en
pourcentage, on cherche une fraction de dénominateur $100$.}{définition 2}
\cas{a)}
\begin{calculs}
p &= \dfrac{12}{40} &&\justif{la partie sur le total}\\
p &= \dfrac{3\times\barre{4}}{10\times\barre{4}} &&\justif{$12=3\times 4$ et $40=10\times 4$}\\
p &= \dfrac{3\fois{10}}{10\fois{10}} &&\justif{on veut $100$ en bas}\\
p &= \dfrac{30}{100} &&
\end{calculs}
\explique{Les $12$ élèves représentent $\res{30\,\%}$ du groupe.}{}
\cas{b)}
\begin{calculs}
p &= \dfrac{9}{20} &&\justif{les admis sur les candidats}\\
p &= \dfrac{9\fois{5}}{20\fois{5}} &&\justif{$20\times 5=100$}\\
p &= \dfrac{45}{100} &&
\end{calculs}
\explique{$\res{45\,\%}$ des candidats sont admis.}{}
\cas{c)}
\begin{calculs}
p &= \dfrac{33}{150} &&\\
p &= \dfrac{11\times\barre{3}}{50\times\barre{3}} &&\justif{$33=11\times 3$ et $150=50\times 3$}\\
p &= \dfrac{11\fois{2}}{50\fois{2}} &&\justif{$50\times 2=100$}\\
p &= \dfrac{22}{100} &&
\end{calculs}
\explique{Les garçons représentent $\res{22\,\%}$ des élèves.}{}
\cas{d)}
\begin{calculs}
p &= \dfrac{80}{500} &&\\
p &= \dfrac{16\times\barre{5}}{100\times\barre{5}} &&\justif{$80=16\times 5$ et $500=100\times 5$}\\
p &= \dfrac{16}{100} &&
\end{calculs}
\explique{Ces clients représentent $\res{16\,\%}$ de l'ensemble.}{}
\rem{Erreur fréquente : diviser le total par la partie ($\dfrac{40}{12}$). Une proportion
est toujours entre $0$ et $1$, donc entre $0\,\%$ et $100\,\%$.}
\end{exo}

\begin{exo}[Deux enquêtes]
\cas{a)}
\begin{calculs}
p &= \dfrac{140}{250} &&\justif{les favorables sur le total}\\
p &= \dfrac{28\times\barre{5}}{50\times\barre{5}} &&\justif{$140=28\times 5$ et $250=50\times 5$}\\
p &= \dfrac{28\fois{2}}{50\fois{2}} &&\justif{$50\times 2=100$}\\
p &= \dfrac{56}{100} &&
\end{calculs}
\explique{$56\,\%$ des élèves sont favorables, et $56\,\%<60\,\%$ : \res{ce n'est pas le
cas}, le projet n'a pas assez de soutien.}{on compare au seuil de $60\,\%$}
\explique{Autre rédaction : $60\,\%$ de $250$ font $150$ ($10\,\%$ font $25$), et
$140<150$.}{on compare au nombre d'élèves qu'il faudrait}
\cas{b)}
\begin{calculs}
p &= \dfrac{340}{400} &&\justif{les satisfaits sur le total}\\
p &= \dfrac{85\times\barre{4}}{100\times\barre{4}} &&\justif{$340=85\times 4$ et $400=100\times 4$}\\
p &= \dfrac{85}{100} &&
\end{calculs}
\explique{$85\,\%$ des clients sont satisfaits : \res{la publicité dit vrai}.}{}
\end{exo}

\begin{exo}[Taux d'évolution]
\explique{Formule du cours : $t=\dfrac{V_{1}-V_{0}}{V_{0}}$, l'arrivée moins le départ,
sur le départ.}{définition 3}
\cas{a)}
\begin{calculs}
t &= \dfrac{180-150}{150} &&\justif{départ $V_{0}=150$, arrivée $V_{1}=180$}\\
t &= \dfrac{30}{150} &&\justif{le prix a pris $30$~€}\\
t &= \dfrac{1}{5} &&\justif{$150=5\times 30$}\\
t &= 0{,}2 &&\justif{$\dfrac{1}{5}=\dfrac{20}{100}$}
\end{calculs}
\explique{$t>0$ : le prix a \res{augmenté de $20\,\%$}.}{$0{,}2=20\,\%$}
\cas{b)}
\begin{calculs}
t &= \dfrac{340-400}{400} &&\justif{départ $400$, arrivée $340$}\\
t &= \dfrac{-60}{400} &&\justif{le stock a perdu $60$~kg}\\
t &= -\dfrac{15}{100} &&\justif{on divise en haut et en bas par $4$}\\
t &= -0{,}15 &&
\end{calculs}
\explique{$t<0$ : le stock a \res{baissé de $15\,\%$}.}{le signe $-$ dit que c'est une
baisse}
\cas{c)}
\begin{calculs}
t &= \dfrac{720-800}{800} &&\\
t &= \dfrac{-80}{800} &&\\
t &= -0{,}1 &&\justif{$\dfrac{80}{800}=\dfrac{1}{10}$}
\end{calculs}
\explique{Le loyer a \res{baissé de $10\,\%$}.}{$0{,}1=10\,\%$}
\cas{d)}
\begin{calculs}
t &= \dfrac{75-60}{60} &&\\
t &= \dfrac{15}{60} &&\\
t &= 0{,}25 &&\justif{$\dfrac{15}{60}=\dfrac{1}{4}$}
\end{calculs}
\explique{Le club a \res{gagné $25\,\%$} d'adhérents.}{}
\rem{Erreur fréquente : diviser par la valeur d'arrivée. En a), $\dfrac{30}{180}$ ne
donne pas $20\,\%$. On divise toujours par la valeur de \emph{départ}.}
\end{exo}

\begin{exo}[Plus ou moins de\ldots{} ?]
\explique{On ne calcule pas le taux : on compare la variation réelle à ce que
représenterait le pourcentage annoncé de la valeur de départ.}{méthode 4, c)}
\cas{a)}
\explique{Le club a perdu $700-620=80$ adhérents.}{la perte réelle}
\explique{$10\,\%$ de $700$ font $70$.}{la perte qu'il faudrait pour $10\,\%$}
\explique{$80>70$ : \res{oui}, le club a perdu plus de $10\,\%$ de ses adhérents.}{}
\explique{Autre rédaction : avec une baisse de $10\,\%$, il resterait $700-70=630$
adhérents ; or $620<630$.}{comme dans les corrigés de l'épreuve}
\cas{b)}
\explique{La forêt a gagné $1\,620-1\,500=120$ arbres.}{le gain réel}
\explique{$10\,\%$ de $1\,500$ font $150$.}{}
\explique{$120<150$ : \res{non}, elle a gagné moins de $10\,\%$ d'arbres.}{}
\cas{c)}
\explique{La ville a perdu $3\,000-2\,880=120$ habitants.}{}
\explique{$10\,\%$ de $3\,000$ font $300$, donc $5\,\%$ en font $150$.}{la moitié de
$10\,\%$}
\explique{$120<150$ : \res{non}, elle a perdu moins de $5\,\%$ de ses habitants.}{}
\end{exo}

\begin{exo}[Au tableur]
\cas{a)}
\begin{calculs}
t &= \dfrac{5\,500-5\,000}{5\,000} &&\justif{départ : 2022 ; arrivée : 2023}\\
t &= \dfrac{500}{5\,000} &&\\
t &= 0{,}1 &&\justif{$\dfrac{500}{5\,000}=\dfrac{1}{10}$}
\end{calculs}
\explique{Le nombre d'entrées a \res{augmenté de $10\,\%$} entre 2022 et 2023.}{}
\cas{b)}
\explique{En C3, on veut le taux de 2022 à 2023 : l'arrivée est en C2, le départ en
B2.}{B3 porte une croix : pas d'année avant 2022}
\explique{La bonne formule est \res{\texttt{=(C2-B2)/B2}}.}{l'arrivée moins le départ, sur
le départ}
\explique{\texttt{=(C2-B2)/C2} divise par l'arrivée ; \texttt{=(B2-C2)/B2} calcule le
départ moins l'arrivée (on obtiendrait $-0{,}1$, une baisse).}{les deux pièges}
\explique{\texttt{=C2-B2/B2} : sans parenthèses, le tableur fait la division d'abord ; il
calcule $5\,500-1=5\,499$.}{la division passe avant la soustraction}
\cas{c)}
\explique{En recopiant vers la droite, chaque lettre avance d'un rang : C devient D, B
devient C.}{rappel de l'énoncé}
\explique{D3 contient \res{\texttt{=(D2-C2)/C2}}.}{le taux de 2023 à 2024}
\rem{Le tableur affiche alors $-0{,}1$ en D3 : de $5\,500$ à $4\,950$, le nombre
d'entrées a baissé de $10\,\%$.}
\end{exo}

\partie{Coefficient multiplicateur}

\begin{exo}[Du taux au coefficient]
\explique{D'où vient le coefficient : augmenter $V$ de $6\,\%$, c'est ajouter
$0{,}06\times V$. On obtient $V+0{,}06\times V=V\times(1+0{,}06)$, soit
$V\times 1{,}06$.}{on garde $100\,\%$ de $V$ et on ajoute $6\,\%$}
\explique{Hausse de $t\,\%$ : $C_{M}=1+\dfrac{t}{100}$ ; baisse de $t\,\%$ :
$C_{M}=1-\dfrac{t}{100}$.}{propriété 1}
\explique{\etiq{a)} $+6\,\%$ : $C_{M}=1+0{,}06=\res{1{,}06}$}{$6\,\%=0{,}06$}
\explique{\etiq{b)} $-6\,\%$ : $C_{M}=1-0{,}06=\res{0{,}94}$}{on garde $94\,\%$ de la
valeur}
\explique{\etiq{c)} $+30\,\%$ : $C_{M}=1+0{,}3=\res{1{,}3}$}{$30\,\%=0{,}30=0{,}3$}
\explique{\etiq{d)} $-30\,\%$ : $C_{M}=1-0{,}3=\res{0{,}7}$}{on garde $70\,\%$}
\explique{\etiq{e)} $+65\,\%$ : $C_{M}=1+0{,}65=\res{1{,}65}$}{}
\explique{\etiq{f)} $-12\,\%$ : $C_{M}=1-0{,}12=\res{0{,}88}$}{on garde $88\,\%$}
\explique{\etiq{g)} $+4\,\%$ : $C_{M}=1+0{,}04=\res{1{,}04}$}{$4\,\%=0{,}04$, et non
$0{,}4$}
\explique{\etiq{h)} $-45\,\%$ : $C_{M}=1-0{,}45=\res{0{,}55}$}{on garde $55\,\%$}
\explique{\etiq{i)} $+200\,\%$ : $C_{M}=1+2=\res{3}$}{$200\,\%=\dfrac{200}{100}=2$ : la
valeur triple}
\rem{Erreurs fréquentes : $+4\,\%$ donne $1{,}04$ et non $1{,}4$ (qui ferait $+40\,\%$) ;
$-6\,\%$ donne $0{,}94$ et non $0{,}6$. Contrôle : une hausse donne un coefficient plus
grand que $1$, une baisse un coefficient plus petit que $1$.}
\end{exo}

\begin{exo}[Du coefficient au taux]
\explique{On écrit le coefficient sous la forme $1+\ldots$ ou $1-\ldots$ : on lit alors
le taux.}{$C_{M}>1$ : hausse ; $C_{M}<1$ : baisse}
\explique{\etiq{a)} $1{,}09=1+0{,}09$ : \res{hausse de $9\,\%$}.}{$C_{M}>1$}
\explique{\etiq{b)} $0{,}91=1-0{,}09$ : \res{baisse de $9\,\%$}.}{$C_{M}<1$}
\explique{\etiq{c)} $\dfrac{4}{5}=0{,}8$ et $0{,}8=1-0{,}2$ : \res{baisse de $20\,\%$}.}{on
écrit d'abord la fraction en décimal}
\explique{\etiq{d)} $0{,}7=1-0{,}3$ : \res{baisse de $30\,\%$}.}{et non de $70\,\%$ : on
garde $70\,\%$, on perd $30\,\%$}
\explique{\etiq{e)} $3=1+2$ : \res{hausse de $200\,\%$}.}{et non de $300\,\%$}
\explique{\etiq{f)} $0{,}98=1-0{,}02$ : \res{baisse de $2\,\%$}.}{}
\rem{Le coefficient dit ce qu'on \emph{garde} ; le taux, ce qu'on \emph{gagne} ou
\emph{perd}. Multiplier par $0{,}7$, c'est garder $70\,\%$, donc perdre $30\,\%$.}
\end{exo}

\begin{exo}[Appliquer une évolution]
\explique{On calcule le coefficient, puis $V_{1}=V_{0}\times C_{M}$ : une seule
multiplication.}{méthode 5}
\cas{a)}
\begin{calculs}
C_{M} &= 1+0{,}15 = 1{,}15 &&\justif{hausse de $15\,\%$}\\
V_{1} &= 300\times 1{,}15 &&\justif{$V_{1}=V_{0}\times C_{M}$}\\
V_{1} &= 345 &&\justif{$300+300\times 0{,}15$, soit $300+45$}
\end{calculs}
\explique{La nouvelle valeur est $\res{345}$~€.}{}
\cas{b)}
\begin{calculs}
C_{M} &= 1-0{,}35 = 0{,}65 &&\justif{baisse de $35\,\%$}\\
V_{1} &= 600\times 0{,}65 &&\\
V_{1} &= 390 &&\justif{$6\times 65=390$}
\end{calculs}
\explique{La nouvelle valeur est $\res{390}$~€.}{}
\cas{c)}
\begin{calculs}
C_{M} &= 1+1{,}2 = 2{,}2 &&\justif{$120\,\%=\dfrac{120}{100}=1{,}2$}\\
V_{1} &= 50\times 2{,}2 &&\\
V_{1} &= 110 &&\justif{$50\times 2=100$ et $50\times 0{,}2=10$}
\end{calculs}
\explique{La nouvelle valeur est $\res{110}$~€ : plus que le double.}{$+120\,\%$, c'est
plus que $+100\,\%$}
\cas{d)}
\begin{calculs}
C_{M} &= 1-0{,}015 = 0{,}985 &&\justif{$1{,}5\,\%=\dfrac{1{,}5}{100}=0{,}015$}\\
V_{1} &= 2\,000\times 0{,}985 &&\\
V_{1} &= 1\,970 &&\justif{$2\,000\times 0{,}015=30$, puis $2\,000-30$}
\end{calculs}
\explique{La commune compte ensuite $\res{1\,970}$ habitants.}{}
\rem{Erreur fréquente en d) : $1{,}5\,\%=0{,}15$. Or $0{,}15=15\,\%$ ; $1{,}5\,\%$ fait
dix fois moins.}
\end{exo}

\begin{exo}[Le bon calcul]
\cas{1.}
\explique{Baisse de $12\,\%$ : $C_{M}=1-0{,}12=0{,}88$, et le nouveau prix est
$500\times 0{,}88$. Réponse \res{b)}.}{on garde $88\,\%$ du prix}
\explique{a) retire $0{,}12$~€ ; c) est une hausse ; d) donne $60$~€, le montant de la
baisse et non le nouveau prix.}{pourquoi les autres sont fausses}
\cas{2.}
\explique{Hausse de $30\,\%$ : $C_{M}=1+0{,}3=1{,}3$, et le nouveau prix est
$150\times 1{,}3$. Réponse \res{c)}.}{}
\explique{a) ajoute $0{,}30$~€ ; b) donne $45$~€, le montant de la hausse ; d) est une
baisse de $30\,\%$.}{}
\end{exo}

\begin{exo}[Retrouver un taux avec le coefficient]
\explique{Comme $V_{1}=V_{0}\times C_{M}$, on trouve $C_{M}=\dfrac{V_{1}}{V_{0}}$ :
l'arrivée sur le départ.}{méthode 5, d)}
\cas{a)}
\begin{calculs}
C_{M} &= \dfrac{100}{80} &&\justif{arrivée sur départ}\\
C_{M} &= 1{,}25 &&\justif{$\dfrac{100}{80}=\dfrac{5}{4}$ ; $1{,}25=1+0{,}25$}
\end{calculs}
\explique{Le prix a \res{augmenté de $25\,\%$}.}{contrôle : $\dfrac{100-80}{80}=\dfrac{20}{80}=0{,}25$}
\cas{b)}
\begin{calculs}
C_{M} &= \dfrac{400}{500} &&\\
C_{M} &= 0{,}8 &&\justif{$\dfrac{4}{5}$ ; $0{,}8=1-0{,}2$}
\end{calculs}
\explique{Le stock a \res{baissé de $20\,\%$}.}{}
\cas{c)}
\begin{calculs}
C_{M} &= \dfrac{60}{20} &&\\
C_{M} &= 3 &&\justif{$3=1+2$}
\end{calculs}
\explique{La production a \res{augmenté de $200\,\%$} : elle a triplé.}{et non de
$300\,\%$}
\end{exo}

\partie{Évolutions successives, évolution réciproque}

\begin{exo}[Taux global]
\explique{Pourquoi on multiplie : après la première évolution, la valeur est
$V_{0}\times C_{M1}$ ; la seconde évolution s'applique à \emph{cette nouvelle valeur},
qui devient $V_{0}\times C_{M1}\times C_{M2}$.}{propriété 2}
\cas{a)}
\begin{calculs}
C_{M} &= 1{,}5\times 1{,}2 &&\justif{$+50\,\%$ : $\times 1{,}5$ ; $+20\,\%$ : $\times 1{,}2$}\\
C_{M} &= 1{,}8 &&\justif{$15\times 12=180$ ; $1{,}8=1+0{,}8$}
\end{calculs}
\explique{Hausse globale de $\res{80\,\%}$.}{et non de $70\,\%$ : avec $100$, on passe à
$150$, puis à $180$}
\cas{b)}
\begin{calculs}
C_{M} &= 0{,}7\times 1{,}3 &&\\
C_{M} &= 0{,}91 &&\justif{$7\times 13=91$ ; $0{,}91=1-0{,}09$}
\end{calculs}
\explique{Baisse globale de $\res{9\,\%}$.}{et non $0\,\%$ : la hausse porte sur une
valeur plus petite}
\cas{c)}
\begin{calculs}
C_{M} &= 2\times 0{,}5 &&\justif{$+100\,\%$ : $\times 2$ ; $-50\,\%$ : $\times 0{,}5$}\\
C_{M} &= 1 &&
\end{calculs}
\explique{Taux global : $\res{0\,\%}$ ; on revient à la valeur de départ.}{doubler puis
prendre la moitié}
\cas{d)}
\begin{calculs}
C_{M} &= 0{,}8\times 0{,}8 &&\\
C_{M} &= 0{,}64 &&\justif{$0{,}64=1-0{,}36$}
\end{calculs}
\explique{Baisse globale de $\res{36\,\%}$.}{et non de $40\,\%$}
\cas{e)}
\begin{calculs}
C_{M} &= 1{,}6\times 0{,}75 &&\justif{$-25\,\%$ : $\times 0{,}75$}\\
C_{M} &= 1{,}2 &&\justif{$\dfrac{3}{4}$ de $1{,}6$ : $1{,}6\div 4=0{,}4$, puis $3\times 0{,}4$}
\end{calculs}
\explique{Hausse globale de $\res{20\,\%}$.}{et non de $35\,\%$}
\rem{Les taux successifs ne s'ajoutent pas : on multiplie les coefficients, puis on lit
le taux sur le coefficient global.}
\end{exo}

\begin{exo}[Le prix de l'électricité]
\begin{calculs}
C_{M} &= 1{,}1\times 1{,}1 &&\justif{deux hausses de $10\,\%$}\\
C_{M} &= 1{,}21 &&\justif{$11\times 11=121$ ; $1{,}21=1+0{,}21$}
\end{calculs}
\explique{Sur l'année, le prix augmente de $21\,\%$ et non de $20\,\%$ : le journaliste
\res{a tort}.}{}
\explique{Avec un prix de $100$~€ : en février, il passe à $110$~€ ; en août, il prend
$10\,\%$ de $110$, soit $11$~€, et passe à $121$~€.}{la seconde hausse porte sur
$110$~€, pas sur $100$~€}
\end{exo}

\begin{exo}[Vrai ou faux ?]
\cas{a)}
\begin{calculs}
C_{M} &= 1{,}5\times 0{,}5 &&\justif{$+50\,\%$, puis $-50\,\%$}\\
C_{M} &= 0{,}75 &&\justif{$0{,}75=1-0{,}25$}
\end{calculs}
\explique{\res{Faux} : on perd $25\,\%$ en tout.}{avec $100$ : $150$, puis la moitié,
$75$}
\cas{b)}
\begin{calculs}
C_{M} &= 0{,}8\times 0{,}7 &&\justif{$-20\,\%$, puis $-30\,\%$}\\
C_{M} &= 0{,}56 &&
\end{calculs}
\explique{Dans l'autre ordre, $0{,}7\times 0{,}8=0{,}56$ aussi : \res{vrai}.}{l'ordre ne
change pas un produit}
\cas{c)}
\begin{calculs}
C_{M} &= 0{,}9\times 0{,}9 &&\\
C_{M} &= 0{,}81 &&\justif{$0{,}81=1-0{,}19$}
\end{calculs}
\explique{\res{Vrai} : la baisse globale est de $19\,\%$, et non de $20\,\%$.}{la seconde
baisse porte sur une valeur plus petite}
\end{exo}

\begin{exo}[La valeur d'un scooter]
\cas{a)}
\begin{calculs}
V_{1} &= 4\,000\times 0{,}75 &&\justif{$-25\,\%$ : $C_{M}=0{,}75$}\\
V_{1} &= 3\,000 &&\justif{les trois quarts de $4\,000$}\\
V_{2} &= 3\,000\times 0{,}8 &&\justif{$-20\,\%$ : $C_{M}=0{,}8$, appliqué à $3\,000$}\\
V_{2} &= 2\,400 &&
\end{calculs}
\explique{Le scooter vaut $\res{3\,000}$~€ au bout d'un an et $\res{2\,400}$~€ au bout
de deux ans.}{}
\cas{b)}
\begin{calculs}
C_{M} &= 0{,}75\times 0{,}8 &&\justif{produit des coefficients}\\
C_{M} &= 0{,}6 &&\justif{$0{,}6=1-0{,}4$}
\end{calculs}
\explique{En deux ans, le scooter a perdu $\res{40\,\%}$ de sa valeur.}{contrôle :
$4\,000\times 0{,}6=2\,400$}
\rem{Erreur fréquente : $25+20=45\,\%$. La seconde baisse porte sur $3\,000$~€ et non
sur $4\,000$~€ : elle retire $600$~€, pas $800$~€.}
\end{exo}

\begin{exo}[Une même évolution répétée]
\cas{a)}
\explique{Trois hausses de $3\,\%$ : $1{,}03\times 1{,}03\times 1{,}03=1{,}03^{3}$.}{la
même évolution trois fois}
\begin{calculs}
C_{M} &= 1{,}03^{3} &&\\
C_{M} &\approx 1{,}09 &&\justif{valeur de l'énoncé ; $1{,}09=1+0{,}09$}
\end{calculs}
\explique{En $3$ ans, la production est multipliée par environ $\res{1{,}09}$ : hausse
d'environ $\res{9\,\%}$.}{« environ » : $1{,}09$ est un arrondi}
\cas{b)}
\explique{Douze hausses de $1\,\%$ : le coefficient est $1{,}01^{12}$.}{$12$ mois}
\begin{calculs}
C_{M} &= 1{,}01^{12} &&\\
C_{M} &\approx 1{,}13 &&\justif{$1{,}13=1+0{,}13$}
\end{calculs}
\explique{Sur un an, les prix augmentent d'environ $\res{13\,\%}$, et non de $12\,\%$.}{chaque
hausse porte sur un prix déjà augmenté}
\end{exo}

\begin{exo}[Se compensent-elles ?]
\explique{Deux évolutions se compensent quand le produit de leurs coefficients vaut
$1$ : on revient alors à la valeur de départ.}{définition 4}
\explique{\etiq{a)} $0{,}8\times 1{,}25=1$ : \res{oui}.}{$-20\,\%$, puis $+25\,\%$}
\explique{\etiq{b)} $1{,}5\times 0{,}5=0{,}75$ : \res{non}.}{il manque $25\,\%$}
\explique{\etiq{c)} $0{,}5\times 2=1$ : \res{oui}.}{la moitié, puis le double}
\explique{\etiq{d)} $4\times 0{,}25=1$ : \res{oui}.}{$+300\,\%$ : $C_{M}=1+3=4$ ;
$-75\,\%$ : $0{,}25$}
\explique{\etiq{e)} $0{,}8\times 1{,}2=0{,}96$ : \res{non}.}{il manque $4\,\%$}
\explique{\etiq{f)} $1{,}1\times 0{,}9=0{,}99$ : \res{non}.}{il manque $1\,\%$}
\rem{Une hausse de $t\,\%$ n'est jamais compensée par une baisse de $t\,\%$ : la baisse
porte sur une valeur plus grande, elle retire donc davantage.}
\end{exo}

\begin{exo}[Quelle évolution compense ?]
\explique{On teste chaque réponse : on cherche celle dont le produit des coefficients
vaut $1$.}{méthode 7}
\cas{1.}
\explique{Baisse de $75\,\%$ : $4\times 0{,}25=1$ ; baisse de $25\,\%$ :
$4\times 0{,}75=3$.}{}
\explique{Une baisse de $400\,\%$ ou de $300\,\%$ est impossible. Réponse \res{b)}.}{on ne
peut pas perdre plus de $100\,\%$}
\cas{2.}
\explique{$1{,}25\times 0{,}75=0{,}9375$ ; \ $1{,}25\times 0{,}25=0{,}3125$ ; \
$1{,}25\times 0{,}8=1$ ; \ $1{,}25\times 0{,}2=0{,}25$.}{coefficients des baisses :
$0{,}75$ ; $0{,}25$ ; $0{,}8$ ; $0{,}2$}
\explique{Il faut une baisse de $20\,\%$. Réponse \res{c)}.}{seul produit égal à $1$}
\cas{3.}
\explique{$0{,}5\times 2=1$ ; \ $0{,}5\times 1{,}5=0{,}75$ ; \ $0{,}5\times 3=1{,}5$ ; \
$0{,}5\times 2{,}5=1{,}25$.}{coefficients des hausses : $2$ ; $1{,}5$ ; $3$ ; $2{,}5$}
\explique{Il faut une hausse de $100\,\%$. Réponse \res{a)}.}{pour revenir de la moitié,
il faut doubler}
\end{exo}

\begin{exo}[Un verger]
\cas{a)}
\begin{calculs}
V_{1} &= 2\,000\times 0{,}8 &&\justif{$-20\,\%$ : on garde $80\,\%$}\\
V_{1} &= 1\,600 &&
\end{calculs}
\explique{Il reste $\res{1\,600}$ arbres.}{}
\cas{b)}
\begin{calculs}
V_{2} &= 1\,600\times 1{,}2 &&\justif{$+20\,\%$ des arbres restants}\\
V_{2} &= 1\,920 &&\justif{$1\,600+320$}
\end{calculs}
\explique{$1\,920<2\,000$ : \res{non}, il lui manque $80$ arbres.}{$20\,\%$ de $1\,600$
($320$) est moins que $20\,\%$ de $2\,000$ ($400$)}
\cas{c)}
\begin{calculs}
V_{2} &= 1\,600\times 1{,}25 &&\justif{$+25\,\%$ : $C_{M}=1{,}25$}\\
V_{2} &= 2\,000 &&\justif{$25\,\%$ de $1\,600$ font $400$}
\end{calculs}
\explique{En replantant $25\,\%$ des arbres restants, il retrouve \res{exactement
$2\,000$ arbres}.}{ou : $0{,}8\times 1{,}25=1$}
\end{exo}

\partie{Fonctions affines : image, antécédent}

\begin{exo}[Affine ou non ?]
\explique{Une fonction est affine quand on peut l'écrire $f(x)=ax+b$ : un nombre fois
$x$, plus un nombre.}{définition 5}
\cas{a)}
\explique{$f(x)=5x+1$ : \res{affine}, avec $a=5$ et $b=1$.}{}
\cas{b)}
\explique{$g(x)=x^{2}-3$ : \res{non affine}.}{$x$ est au carré}
\cas{c)}
\explique{$h(x)=2-3x$ s'écrit $h(x)=-3x+2$ : \res{affine}, avec $a=-3$ et $b=2$.}{$a$ est
le nombre devant $x$, avec son signe}
\cas{d)}
\explique{$k(x)=\dfrac{2}{x}$ : \res{non affine}.}{$x$ est au dénominateur}
\cas{e)}
\explique{$m(x)=1{,}5x$ s'écrit $m(x)=1{,}5x+0$ : \res{affine}, avec $a=1{,}5$ et $b=0$.}{$b=0$ :
fonction linéaire}
\cas{f)}
\explique{$n(x)=-4$ s'écrit $n(x)=0\times x-4$ : \res{affine}, avec $a=0$ et
$b=-4$.}{$a=0$ : fonction constante}
\rem{Erreur fréquente en c) : lire $a=2$ parce que $2$ est écrit en premier. Le
coefficient directeur est celui de $x$, où qu'il soit placé.}
\end{exo}

\begin{exo}[Images]
\explique{Pour calculer une image, on remplace $x$ par le nombre, entre parenthèses
s'il est négatif.}{méthode 8}
\cas{a)}
\begin{calculs}
f(0) &= 3\times 0-4 &&\justif{on remplace $x$ par $0$}\\
f(0) &= \res{-4} &&\justif{$f(0)=b$, toujours}\\
f(2) &= 3\times 2-4 &&\\
f(2) &= \res{2} &&\justif{$6-4$}\\
f(-1) &= 3\times(-1)-4 &&\justif{parenthèses autour de $-1$}\\
f(-1) &= \res{-7} &&\justif{$-3-4$}\\
f(0{,}5) &= 3\times 0{,}5-4 &&\\
f(0{,}5) &= \res{-2{,}5} &&\justif{$1{,}5-4$}
\end{calculs}
\cas{b)}
\begin{calculs}
g(0) &= -4\times 0+6 &&\\
g(0) &= \res{6} &&\\
g(2) &= -4\times 2+6 &&\\
g(2) &= \res{-2} &&\justif{$-8+6$}\\
g(-3) &= -4\times(-3)+6 &&\justif{parenthèses autour de $-3$}\\
g(-3) &= \res{18} &&\justif{$-4\times(-3)=12$ : moins par moins donne plus}\\
g(1{,}5) &= -4\times 1{,}5+6 &&\\
g(1{,}5) &= \res{0} &&\justif{$-6+6$}
\end{calculs}
\rem{Erreur fréquente : écrire $-4\times -3$ sans parenthèses, puis se tromper de
signe. $g(-3)=12+6$.}
\end{exo}

\begin{exo}[Antécédents]
\explique{Chercher l'antécédent de $15$, c'est trouver $x$ tel que $f(x)=15$ : on
résout une équation.}{méthode 9}
\cas{a)}
\begin{calculs}
3x+9 &= 15 &&\justif{on écrit $f(x)=15$}\\
3x &= 6 &&\justif{on retire $9$ des deux côtés}\\
x &= 2 &&\justif{on divise par $3$}
\end{calculs}
\explique{L'antécédent de $15$ est $\res{2}$.}{contrôle : $f(2)=6+9=15$}
\begin{calculs}
3x+9 &= 0 &&\justif{on écrit $f(x)=0$}\\
3x &= -9 &&\justif{on retire $9$}\\
x &= -3 &&\justif{on divise par $3$}
\end{calculs}
\explique{L'antécédent de $0$ est $\res{-3}$.}{contrôle : $f(-3)=-9+9=0$}
\begin{calculs}
3x+9 &= -6 &&\justif{on écrit $f(x)=-6$}\\
3x &= -15 &&\justif{$-6-9=-15$}\\
x &= -5 &&\justif{on divise par $3$}
\end{calculs}
\explique{L'antécédent de $-6$ est $\res{-5}$.}{contrôle : $f(-5)=-15+9=-6$}
\cas{b)}
\begin{calculs}
-2x+8 &= 0 &&\justif{on écrit $g(x)=0$}\\
-2x &= -8 &&\justif{on retire $8$}\\
x &= 4 &&\justif{on divise par $-2$ : $\dfrac{-8}{-2}=4$}
\end{calculs}
\explique{L'antécédent de $0$ est $\res{4}$.}{contrôle : $g(4)=-8+8=0$}
\begin{calculs}
-2x+8 &= 14 &&\justif{on écrit $g(x)=14$}\\
-2x &= 6 &&\justif{on retire $8$}\\
x &= -3 &&\justif{on divise par $-2$ : $\dfrac{6}{-2}=-3$}
\end{calculs}
\explique{L'antécédent de $14$ est $\res{-3}$.}{contrôle : $g(-3)=6+8=14$}
\cas{c)}
\begin{calculs}
0{,}5x-2 &= 3 &&\justif{on recopie l'équation}\\
0{,}5x &= 5 &&\justif{on ajoute $2$}\\
x &= 10 &&\justif{diviser par $0{,}5$, c'est multiplier par $2$}
\end{calculs}
\explique{La solution est $\res{10}$.}{contrôle : $h(10)=5-2=3$}
\rem{Erreur fréquente en b) : oublier le signe de $-2$ et trouver $x=3$. Le contrôle
le montre : $g(3)=-6+8=2$, et non $14$.}
\end{exo}

\begin{exo}[Deux tarifs]
\cas{a)}
\explique{Le prix est la partie fixe, $40$~€, plus $1{,}5$~€ pour chacun des $x$
kilomètres : $T(x)=\res{1{,}5x+40}$.}{$b=40$ : le fixe ; $a=1{,}5$ : le prix d'un km}
\begin{calculs}
1{,}5x+40 &= 70 &&\justif{on écrit $T(x)=70$}\\
1{,}5x &= 30 &&\justif{on retire $40$}\\
x &= 20 &&\justif{$\dfrac{30}{1{,}5}=\dfrac{300}{15}$}
\end{calculs}
\explique{Le dépanneur a parcouru $\res{20}$~km.}{contrôle : $T(20)=30+40=70$}
\cas{b)}
\explique{$S(x)=\res{15x+200}$.}{$200$~€ fixes, plus $15$~€ par invité}
\begin{calculs}
S(40) &= 15\times 40+200 &&\justif{on remplace $x$ par $40$}\\
S(40) &= 800 &&\justif{$600+200$}
\end{calculs}
\explique{Pour $40$ invités, la salle coûte $\res{800}$~€.}{}
\begin{calculs}
15x+200 &= 650 &&\justif{on écrit $S(x)=650$}\\
15x &= 450 &&\justif{on retire $200$}\\
x &= 30 &&\justif{$\dfrac{450}{15}=30$}
\end{calculs}
\explique{Pour $650$~€, on peut inviter $\res{30}$ personnes.}{contrôle : $450+200=650$}
\end{exo}

\partie{Trouver une fonction affine}

\begin{exo}[Mettre en équation]
\explique{Dans $ax+b$, $b$ est la valeur fixe (celle pour $x=0$) et $a$ ce que coûte,
rapporte ou fait perdre une unité de $x$.}{propriété 3}
\cas{a)}
\explique{$b=8$ (le fixe) et $a=6$ (le prix d'une heure) : $L(x)=\res{6x+8}$.}{contrôle :
$L(1)=14$~€ pour une heure}
\cas{b)}
\explique{$5$~L pour $100$~km, c'est $\dfrac{5}{100}=0{,}05$~L pour $1$~km.}{on se ramène à
une unité de $x$ : $1$~km}
\explique{Après $x$~km, le réservoir a perdu $0{,}05x$~L : $R(x)=60-0{,}05x$, soit
$R(x)=\res{-0{,}05x+60}$.}{$a<0$ : la quantité diminue}
\explique{Contrôle : $R(100)=60-5=55$ ; après $100$~km, il reste bien $55$~L.}{}
\cas{c)}
\explique{$4\,\%$ des ventes, c'est $0{,}04\times x$ : $S(x)=\res{0{,}04x+1\,400}$.}{$4\,\%
=0{,}04$ ; le fixe est $1\,400$~€}
\rem{Erreur fréquente en b) : $R(x)=60-5x$. Le contrôle le montre : on aurait
$R(100)=60-500$, un nombre négatif.}
\end{exo}

\begin{exo}[Coefficient directeur]
\explique{$a=\dfrac{y_{B}-y_{A}}{x_{B}-x_{A}}$ : différence des ordonnées sur différence
des abscisses, dans le même ordre.}{propriété 3}
\cas{a)}
\begin{calculs}
a &= \dfrac{11-2}{4-1} &&\justif{$y_{B}-y_{A}$ en haut, $x_{B}-x_{A}$ en bas}\\
a &= \dfrac{9}{3} &&\\
a &= \res{3} &&
\end{calculs}
\cas{b)}
\begin{calculs}
a &= \dfrac{2-6}{2-0} &&\\
a &= \dfrac{-4}{2} &&\\
a &= \res{-2} &&\justif{la droite descend}
\end{calculs}
\cas{c)}
\begin{calculs}
a &= \dfrac{7-3}{-4-(-2)} &&\justif{parenthèses autour de $-2$}\\
a &= \dfrac{4}{-2} &&\justif{$-4-(-2)=-4+2=-2$}\\
a &= \res{-2} &&
\end{calculs}
\cas{d)}
\begin{calculs}
a &= \dfrac{6-(-2)}{1-(-3)} &&\\
a &= \dfrac{8}{4} &&\justif{$6+2=8$ et $1+3=4$}\\
a &= \res{2} &&
\end{calculs}
\rem{On peut aussi commencer par $A$ en haut \emph{et} en bas : $\dfrac{2-11}{1-4}=
\dfrac{-9}{-3}=3$. Mélanger les deux ordres change le signe du résultat.}
\end{exo}

\begin{exo}[Déterminer une fonction affine]
\explique{On calcule $a$ avec les deux points, puis $b$ en écrivant que l'un des points
est sur la droite ; on contrôle avec l'autre.}{méthode 10}
\cas{a)}
\begin{calculs}
a &= \dfrac{11-5}{3-1} &&\justif{$\dfrac{y_{B}-y_{A}}{x_{B}-x_{A}}$}\\
a &= 3 &&\justif{$\dfrac{6}{2}$}\\
3\times 1+b &= 5 &&\justif{le point $A(1\,;5)$ est sur la droite : $f(1)=5$}\\
b &= 2 &&\justif{on retire $3$}
\end{calculs}
\explique{$f(x)=\res{3x+2}$.}{contrôle avec $B$ : $f(3)=9+2=11$}
\cas{b)}
\explique{$A(0\,;-4)$ a pour abscisse $0$ : la droite coupe l'axe des ordonnées en $-4$,
donc $b=-4$.}{$b=f(0)$ : rien à calculer}
\begin{calculs}
a &= \dfrac{2-(-4)}{2-0} &&\\
a &= 3 &&\justif{$\dfrac{6}{2}$}
\end{calculs}
\explique{$f(x)=\res{3x-4}$.}{contrôle avec $B$ : $f(2)=6-4=2$}
\cas{c)}
\begin{calculs}
a &= \dfrac{3-3}{2-(-2)} &&\\
a &= 0 &&\justif{$\dfrac{0}{4}$}
\end{calculs}
\explique{$a=0$ : $f$ est constante, égale à l'ordonnée commune des deux points :
$f(x)=\res{3}$.}{droite horizontale}
\cas{d)}
\begin{calculs}
a &= \dfrac{4-(-5)}{2-(-4)} &&\justif{parenthèses autour des négatifs}\\
a &= \dfrac{9}{6} &&\justif{$4+5=9$ et $2+4=6$}\\
a &= 1{,}5 &&\justif{$\dfrac{9}{6}=\dfrac{3}{2}$}\\
1{,}5\times 2+b &= 4 &&\justif{avec le point $B(2\,;4)$ : calcul plus simple}\\
3+b &= 4 &&\\
b &= 1 &&
\end{calculs}
\explique{$f(x)=\res{1{,}5x+1}$.}{contrôle avec $A$ : $f(-4)=-6+1=-5$}
\end{exo}

\begin{exo}[Le dioxyde de carbone dans l'air]
\cas{a)}
\explique{$x$ compte les années depuis 2000 : en 2000, $x=0$ ; en 2010, $x=10$. Donc
$f(0)=\res{370}$ et $f(10)=\res{390}$.}{$x$ n'est pas l'année elle-même}
\explique{$f(0)=370$ donne directement $b=370$.}{$b=f(0)$ : la valeur de départ}
\begin{calculs}
a &= \dfrac{390-370}{10-0} &&\justif{points $(0\,;370)$ et $(10\,;390)$}\\
a &= 2 &&\justif{$\dfrac{20}{10}$}
\end{calculs}
\explique{$f(x)=\res{2x+370}$.}{contrôle : $f(10)=20+370=390$}
\cas{b)}
\explique{Quand $x$ augmente de $1$, $f(x)$ augmente de $2$ : selon ce modèle, la
concentration \res{augmente de $2$~ppm par an}.}{$a$ : la variation pour une année}
\cas{c)}
\begin{calculs}
f(30) &= 2\times 30+370 &&\justif{2030 : $x=2030-2000=30$}\\
f(30) &= 430 &&
\end{calculs}
\explique{En 2030, la concentration serait d'environ $\res{430}$~ppm.}{}
\cas{d)}
\begin{calculs}
2x+370 &= 450 &&\justif{on écrit $f(x)=450$}\\
2x &= 80 &&\justif{on retire $370$}\\
x &= 40 &&
\end{calculs}
\explique{$x=40$ correspond à $2000+40$ : la concentration atteindrait $450$~ppm en
\res{2040}.}{on revient à l'année}
\rem{Erreur fréquente : calculer $f(2030)$. On remplace $x$ par le nombre d'années
écoulées depuis 2000, pas par l'année.}
\end{exo}

\partie{Tracer une fonction affine}

\begin{exo}[Sens de variation]
\explique{Seul le signe de $a$ compte : $a>0$, croissante ; $a<0$, décroissante ; $a=0$,
constante.}{propriété 4}
\explique{\etiq{a)} $f(x)=-x+4$ : $a=-1<0$, $f$ est \res{décroissante}.}{$4-x$ s'écrit
$-1\times x+4$}
\explique{\etiq{b)} $a=3>0$ : $g$ est \res{croissante}.}{}
\explique{\etiq{c)} $k(x)=-2$ : $a=0$, $k$ est \res{constante}.}{$-2$ est $b$, pas $a$}
\explique{\etiq{d)} $a=-5<0$ : $p$ est \res{décroissante}.}{}
\explique{\etiq{e)} $a=-0{,}2<0$ : $h$ est \res{décroissante}.}{même si $b=3$ est positif}
\explique{\etiq{f)} $a=0{,}3>0$ : $m$ est \res{croissante}.}{même si $b=-1$ est négatif}
\rem{Le signe de $b$ ne joue aucun rôle dans le sens de variation.}
\end{exo}

\begin{exo}[Tracer des droites]
\cotefig[0.42]{%
\explique{Pour chaque droite, on calcule deux images, on place les deux points
(croix), puis on trace.}{deux points suffisent}
\explique{$f(0)=1$ et $f(2)=3$ : points $(0\,;1)$ et $(2\,;3)$.}{}
\explique{$g(0)=2$ et $g(1)=-1$ : points $(0\,;2)$ et $(1\,;-1)$.}{$-3\times 1+2$}
\explique{$h(0)=-2$ et $h(4)=0$ : points $(0\,;-2)$ et $(4\,;0)$.}{abscisse paire : pas
de demi-unité}
\explique{Contrôle : $f$ et $h$ montent ($a>0$) ; $g$ descend, et vite ($a=-3$).}{pour $1$
vers la droite, on descend de $3$}}{%
\begin{tikzpicture}[x=0.5cm,y=0.5cm]
  \repere{-3}{-4}{5}{4}
  \gradx{-2/-2,-1/-1,1/1,2/2,3/3,4/4}
  \grady{-3/-3,-2/-2,-1/-1,1/1,2/2,3/3}
  \begin{scope}
    \clip (-3,-4) rectangle (5,4);
    \draw[coulCourbe,thick] (-3,-2) -- (3,4);
    \draw[coulCourbeB,thick] (-0.667,4) -- (2,-4);
    \draw[coulCourbeC,thick] (-3,-3.5) -- (5,0.5);
  \end{scope}
  \foreach \p in {(0,1),(2,3)} {\path[coulCourbe] \p node[croix]{};}
  \foreach \p in {(0,2),(1,-1)} {\path[coulCourbeB] \p node[croix]{};}
  \foreach \p in {(0,-2),(4,0)} {\path[coulCourbeC] \p node[croix]{};}
  \node[coulCourbe,font=\small,left] at (2.6,3.7) {$f$};
  \node[coulCourbeB,font=\small,right] at (1.7,-3.4) {$g$};
  \node[coulCourbeC,font=\small,above] at (4.4,0.3) {$h$};
\end{tikzpicture}}
\end{exo}

\begin{exo}[Location de paddles]
\cas{a)}
\explique{$L(x)=\res{4x+4}$.}{$4$~€ fixes, plus $4$~€ par heure}
\cas{b)}
\begin{calculs}
L(0) &= 4\times 0+4 &&\\
L(0) &= \res{4} &&\justif{on paie le fixe même sans utiliser le paddle}\\
L(8) &= 4\times 8+4 &&\\
L(8) &= \res{36} &&\justif{$32+4$}
\end{calculs}
\cas{c) et d)}
\cotefig[0.5]{%
\explique{c) On place $(0\,;4)$ et $(8\,;36)$, calculés en b), puis on trace le segment
qui les relie.}{$x$ va de $0$ à $8$ : on ne prolonge pas}
\explique{d) On part de $24$ sur l'axe des ordonnées, on va horizontalement jusqu'à la
droite, puis on descend sur l'axe des abscisses : on lit $5$.}{lecture en pointillés}
\begin{calculs}
4x+4 &= 24 &&\justif{on écrit $L(x)=24$}\\
4x &= 20 &&\justif{on retire $4$}\\
x &= 5 &&\justif{on divise par $4$}
\end{calculs}
\explique{Pour $24$~€, on loue le paddle $\res{5}$ heures.}{le calcul confirme la
lecture}}{%
\begin{tikzpicture}[x=1cm,y=0.2cm]
  \draw[grille,xstep=1,ystep=5] (0,0) grid (8,40);
  \draw[->,>=Stealth,black!70] (0,0) -- (8.6,0) node[below,font=\scriptsize] {heures};
  \draw[->,>=Stealth,black!70] (0,0) -- (0,43) node[left,font=\scriptsize] {€};
  \foreach \i in {1,...,8} {\node[below,font=\scriptsize,inner sep=1.5pt] at (\i,0) {$\i$};}
  \foreach \v in {5,10,...,40} {\node[left,font=\scriptsize,inner sep=1.5pt] at (0,\v) {$\v$};}
  \node[below left,font=\scriptsize,inner sep=1.5pt] at (0,0) {$0$};
  \draw[coulCourbe,thick] (0,4) -- (8,36);
  \draw[lecture] (0,24) -- (5,24) -- (5,0);
  \foreach \p in {(0,4),(8,36),(5,24)} {\path \p node[croix]{};}
\end{tikzpicture}}
\rem{Repère non orthonormé : les deux axes portent des grandeurs différentes (heures et
euros), à l'échelle de l'énoncé.}
\end{exo}

\begin{exo}[Proportionnalité et fonctions linéaires]
\cas{a)}
\explique{$10$~g pour $100$~mL, c'est $\dfrac{10}{100}=0{,}1$~g pour $1$~mL :
$f(x)=\res{0{,}1x}$.}{la quantité pour une unité : le coefficient $a$}
\begin{calculs}
f(330) &= 0{,}1\times 330 &&\\
f(330) &= 33 &&\justif{multiplier par $0{,}1$, c'est diviser par $10$}
\end{calculs}
\explique{$330$~mL de boisson contiennent $\res{33}$~g de sucre.}{}
\cas{b)}
\explique{On note $c(x)=2{,}5x$ la longueur en centimètres de $x$ pouces.}{proportionnalité :
fonction linéaire}
\begin{calculs}
c(40) &= 2{,}5\times 40 &&\\
c(40) &= 100 &&\justif{$2\times 40+0{,}5\times 40$, soit $80+20$}
\end{calculs}
\explique{$40$ pouces font environ $\res{100}$~cm.}{}
\begin{calculs}
2{,}5x &= 75 &&\justif{on cherche l'antécédent de $75$}\\
x &= 30 &&\justif{$\dfrac{75}{2{,}5}=\dfrac{750}{25}$}
\end{calculs}
\explique{$75$~cm font environ $\res{30}$ pouces.}{}
\cas{c)}
\explique{Une fonction linéaire s'écrit $ax$ : pas de nombre ajouté ($b=0$).}{propriété 5}
\explique{$f(x)=-2x$ et $h(x)=\dfrac{x}{4}=0{,}25x$ sont \res{linéaires}.}{$b=0$}
\explique{$g(x)=2x+3$ n'est pas linéaire ($b=3$) ; $k(x)=-1$ non plus (constante,
$b=-1$).}{affines, mais pas linéaires}
\cas{d)}
\begin{calculs}
\ell(5) &= 15 &&\justif{le point $(5\,;15)$ est sur la droite}\\
a\times 5 &= 15 &&\justif{$\ell(x)=ax$}\\
a &= 3 &&\justif{on divise par $5$}
\end{calculs}
\explique{$\ell(x)=\res{3x}$.}{contrôle : $\ell(5)=15$}
\end{exo}

\partie{Lire l'équation d'une droite}

\begin{exo}[Lire une équation réduite]
\explique{$b$ : l'ordonnée du point où la droite coupe l'axe des ordonnées. $a$ : en
partant de ce point, on avance jusqu'à un nœud du quadrillage sur la droite, et
$a=\dfrac{\text{déplacement vertical}}{\text{déplacement horizontal}}$.}{propriété 6 ;
une descente compte en négatif}
\cas{$(d_{1})$}
\cotefig{%
\explique{$b=-1$}{la droite coupe l'axe des ordonnées en $(0\,;-1)$}
\explique{De $(0\,;-1)$ à $(1\,;2)$ : on avance de $1$ et on monte de $3$, donc
$a=\dfrac{3}{1}=3$.}{vertical sur horizontal}
\explique{$(d_{1})$ : $\res{y=3x-1}$ ; la fonction est \res{croissante}.}{$a>0$}}{%
\begin{tikzpicture}[x=0.4cm,y=0.4cm]
  \reperefiche
  \begin{scope}
    \clip (-4,-4) rectangle (5,5);
    \draw[coulCourbe,thick] (-1,-4) -- (2,5);
  \end{scope}
  \draw[pas,coulCourbe] (0,-1) -- node[below,font=\footnotesize]{$1$} (1,-1)
    -- node[right,font=\footnotesize]{$+3$} (1,2);
  \foreach \p in {(0,-1),(1,2)} {\path \p node[croix]{};}
\end{tikzpicture}}
\cas{$(d_{2})$}
\cotefig{%
\explique{$b=2$}{point $(0\,;2)$}
\explique{De $(0\,;2)$ à $(2\,;0)$ : on avance de $2$ et on descend de $2$, donc
$a=\dfrac{-2}{2}=-1$.}{on descend : déplacement négatif}
\explique{$(d_{2})$ : $\res{y=-x+2}$ ; la fonction est \res{décroissante}.}{$a<0$ ;
$-1\times x$ s'écrit $-x$}}{%
\begin{tikzpicture}[x=0.4cm,y=0.4cm]
  \reperefiche
  \begin{scope}
    \clip (-4,-4) rectangle (5,5);
    \draw[coulCourbeB,thick] (-3,5) -- (5,-3);
  \end{scope}
  \draw[pas,coulCourbeB] (0,2) -- node[above,font=\footnotesize]{$2$} (2,2)
    -- node[right,font=\footnotesize]{$-2$} (2,0);
  \foreach \p in {(0,2),(2,0)} {\path \p node[croix]{};}
\end{tikzpicture}}
\cas{$(d_{3})$}
\cotefig{%
\explique{$b=1$}{point $(0\,;1)$}
\explique{La droite ne repasse sur un nœud qu'en $(3\,;2)$ : on avance de $3$ et on monte
de $1$, donc $a=\dfrac{1}{3}$.}{on cherche un nœud du quadrillage sur la droite}
\explique{$(d_{3})$ : $\res{y=\dfrac{1}{3}x+1}$ ; la fonction est
\res{croissante}.}{$a>0$}}{%
\begin{tikzpicture}[x=0.4cm,y=0.4cm]
  \reperefiche
  \draw[coulCourbeC,thick] (-4,-0.333) -- (5,2.667);
  \draw[pas,coulCourbeC] (0,1) -- node[below,font=\footnotesize]{$3$} (3,1)
    -- node[right,font=\footnotesize]{$+1$} (3,2);
  \foreach \p in {(0,1),(3,2)} {\path \p node[croix]{};}
\end{tikzpicture}}
\cas{$(d_{4})$}
\cotefig{%
\explique{La droite est horizontale : $a=0$.}{on avance sans monter ni descendre}
\explique{Elle coupe l'axe des ordonnées en $-2$ : $b=-2$.}{}
\explique{$(d_{4})$ : $\res{y=-2}$ ; la fonction est \res{constante}.}{$y=0\times x-2$}}{%
\begin{tikzpicture}[x=0.4cm,y=0.4cm]
  \reperefiche
  \draw[black!70,thick] (-4,-2) -- (5,-2);
  \path (0,-2) node[croix]{};
\end{tikzpicture}}
\rem{Erreur fréquente pour $(d_{3})$ : $a=3$. On a avancé de $3$ pour monter de $1$ :
$a=\dfrac{1}{3}$, la droite monte doucement.}
\end{exo}

\begin{exo}[Points d'une droite]
\cas{a)}
\explique{Un point est sur $(d)$ quand son ordonnée est égale à $-2\times x+5$, calculé
avec son abscisse.}{propriété 4 : $M(x\,;y)$ est sur la droite quand $y=ax+b$}
\explique{$A(1\,;3)$ : $-2\times 1+5=3$, l'ordonnée de $A$ : \res{$A$ est sur
$(d)$}.}{}
\explique{$B(3\,;1)$ : $-2\times 3+5=-1$, et $-1\ne 1$ : \res{$B$ n'est pas sur
$(d)$}.}{}
\explique{$C(-2\,;9)$ : $-2\times(-2)+5=9$, l'ordonnée de $C$ : \res{$C$ est sur
$(d)$}.}{$4+5$}
\cas{b)}
\begin{calculs}
y &= -2\times 4+5 &&\justif{on remplace $x$ par $4$}\\
y &= -3 &&\justif{$-8+5$}
\end{calculs}
\explique{Le point de $(d)$ d'abscisse $4$ a pour ordonnée $\res{-3}$.}{}
\cas{c)}
\begin{calculs}
-2x+5 &= -7 &&\justif{on remplace $y$ par $-7$}\\
-2x &= -12 &&\justif{on retire $5$}\\
x &= 6 &&\justif{on divise par $-2$}
\end{calculs}
\explique{Le point de $(d)$ d'ordonnée $-7$ a pour abscisse $\res{6}$.}{contrôle :
$-12+5=-7$}
\end{exo}

\begin{exo}[Équation réduite sans graphique]
\cas{1.}
\explique{La droite coupe l'axe des ordonnées en $2$ : $b=2$.}{}
\explique{$a=\dfrac{\text{vertical}}{\text{horizontal}}=\dfrac{1}{3}$. Réponse
\res{b)} : $y=\dfrac{1}{3}x+2$.}{a) inverse les deux déplacements ; d) a $b=-2$}
\cas{2.}
\explique{Le point $(0\,;-3)$ donne $b=-3$.}{abscisse $0$}
\begin{calculs}
a &= \dfrac{1-(-3)}{2-0} &&\\
a &= 2 &&\justif{$\dfrac{4}{2}$}
\end{calculs}
\explique{Réponse \res{c)} : $y=2x-3$.}{contrôle : $2\times 2-3=1$}
\cas{3.}
\explique{$a=4>0$ : la droite \res{monte}, réponse \res{a)}.}{}
\end{exo}

\partie{Signe d'une fonction affine}

\begin{exo}[Tableaux de signes]
\explique{On résout $ax+b=0$ pour placer le $0$ ; à droite du $0$, on écrit le signe de
$a$, à gauche le signe contraire.}{propriété 7}
\cas{a)}
\cotefig[0.42]{%
\begin{calculs}
3x-9 &= 0 &&\justif{où placer le $0$}\\
3x &= 9 &&\justif{on ajoute $9$}\\
x &= 3 &&
\end{calculs}
\explique{$a=3>0$ : $+$ à droite de $3$.}{la droite monte}
\explique{Contrôle : $f(0)=-9$, négatif, et $0$ est à gauche de $3$.}{on teste une
valeur}}{\signes{f(x)}{3}{-}{+}}
\cas{b)}
\cotefig[0.42]{%
\begin{calculs}
-2x+8 &= 0 &&\\
-2x &= -8 &&\justif{on retire $8$}\\
x &= 4 &&\justif{on divise par $-2$}
\end{calculs}
\explique{$a=-2<0$ : $-$ à droite de $4$.}{la droite descend}
\explique{Contrôle : $g(0)=8$, positif.}{}}{\signes{g(x)}{4}{+}{-}}
\cas{c)}
\cotefig[0.42]{%
\begin{calculs}
4x+12 &= 0 &&\\
4x &= -12 &&\justif{on retire $12$}\\
x &= -3 &&
\end{calculs}
\explique{$a=4>0$ : $+$ à droite de $-3$.}{}
\explique{Contrôle : $h(0)=12$, positif, et $0$ est à droite de $-3$.}{}}{\signes{h(x)}{-3}{-}{+}}
\cas{d)}
\cotefig[0.42]{%
\begin{calculs}
-0{,}5x+3 &= 0 &&\\
-0{,}5x &= -3 &&\justif{on retire $3$}\\
x &= 6 &&\justif{$\dfrac{-3}{-0{,}5}=\dfrac{30}{5}$}
\end{calculs}
\explique{$a=-0{,}5<0$ : $-$ à droite de $6$.}{}
\explique{Contrôle : $k(0)=3$, positif.}{}}{\signes{k(x)}{6}{+}{-}}
\rem{Erreur fréquente : recopier le signe de $b$ (en a), $-9$) à droite du $0$. C'est le
signe de $a$ qui va à droite.}
\end{exo}

\begin{exo}[Positif ou négatif ?]
\cas{1.}
\cotefig[0.42]{%
\begin{calculs}
-5x+2 &= 0 &&\\
-5x &= -2 &&\\
x &= 0{,}4 &&\justif{$\dfrac{2}{5}=0{,}4$}
\end{calculs}
\explique{$a=-5<0$ : $f(x)$ est positif avant $0{,}4$, négatif après.}{}
\explique{$f(x)$ est positif ou nul pour $x\le 0{,}4$ : réponse \res{a)}.}{contrôle :
$f(0)=2$}}{\signes{f(x)}{0{,}4}{+}{-}}
\cas{2.}
\cotefig[0.42]{%
\begin{calculs}
3x+12 &= 0 &&\\
3x &= -12 &&\\
x &= -4 &&
\end{calculs}
\explique{$a=3>0$ : $3x+12$ est négatif avant $-4$.}{}
\explique{Strictement négatif pour $x<-4$ : réponse \res{c)}.}{contrôle avec $-5$ :
$-15+12=-3$}}{\signes{3x+12}{-4}{-}{+}}
\end{exo}

\begin{exo}[Un bénéfice]
\cas{a)}
\begin{calculs}
B(0) &= 5\times 0-150 &&\\
B(0) &= \res{-150} &&
\end{calculs}
\explique{Sans aucun repas vendu, le food-truck perd $150$~€ : ce sont ses frais fixes
de la journée.}{bénéfice négatif : une perte}
\cas{b)}
\begin{calculs}
5x-150 &= 0 &&\\
5x &= 150 &&\justif{on ajoute $150$}\\
x &= \res{30} &&\justif{on divise par $5$}
\end{calculs}
\cas{c)}
\cotefig[0.42]{%
\explique{$a=5>0$ : $B(x)$ est négatif avant $30$ et positif après.}{chaque repas
rapporte $5$~€}
\explique{Le tableau commence à $0$ et finit à $80$.}{$x$ va de $0$ à $80$}}{%
\begin{tikzpicture}
\tkzTabInit[lgt=1.4,espcl=1.5]{$x$/0.7,$B(x)$/0.7}{$0$,$30$,$80$}
\tkzTabLine{,-,z,+,}
\end{tikzpicture}}
\cas{d)}
\explique{Le bénéfice est nul pour $30$ repas et positif au-delà : il est positif à
partir de $\res{31}$ repas vendus.}{$B(31)=155-150=5$}
\end{exo}

\partie{Comparer deux fonctions affines}

\begin{exo}[Équations]
\explique{On regroupe les $x$ d'un côté et les nombres de l'autre, en faisant la même
opération des deux côtés, puis on divise par le nombre devant $x$.}{rappel de la
partie X}
\cas{a)}
\begin{calculs}
6x+5 &= 4x+11 &&\justif{on recopie l'équation}\\
2x+5 &= 11 &&\justif{on retire $4x$ des deux côtés}\\
2x &= 6 &&\justif{on retire $5$}\\
x &= \res{3} &&\justif{on divise par $2$}
\end{calculs}
\explique{Contrôle : $6\times 3+5=23$ et $4\times 3+11=23$.}{}
\cas{b)}
\begin{calculs}
3x-2 &= x+8 &&\\
2x-2 &= 8 &&\justif{on retire $x$}\\
2x &= 10 &&\justif{on ajoute $2$}\\
x &= \res{5} &&
\end{calculs}
\explique{Contrôle : $15-2=13$ et $5+8=13$.}{}
\cas{c)}
\begin{calculs}
2x+3 &= 5x-9 &&\\
3 &= 3x-9 &&\justif{on retire $2x$ : les $x$ restent à droite, en positif}\\
12 &= 3x &&\justif{on ajoute $9$}\\
x &= \res{4} &&\justif{on divise par $3$}
\end{calculs}
\explique{Contrôle : $8+3=11$ et $20-9=11$.}{}
\cas{d)}
\begin{calculs}
4x+1 &= -2x+4 &&\\
6x+1 &= 4 &&\justif{on ajoute $2x$ : $4x+2x=6x$}\\
6x &= 3 &&\justif{on retire $1$}\\
x &= \res{0{,}5} &&\justif{$\dfrac{3}{6}=\dfrac{1}{2}$}
\end{calculs}
\explique{Contrôle : $2+1=3$ et $-1+4=3$.}{}
\rem{Erreur fréquente en d) : retirer $-2x$ en écrivant $4x-2x$. Pour faire
disparaître $-2x$, on \emph{ajoute} $2x$.}
\end{exo}

\begin{exo}[Deux offres d'emploi]
\cas{a)}
\explique{Société A : $10\,\%$ des ventes, c'est $0{,}1x$ : $f(x)=\res{0{,}1x+1\,200}$.}{le
fixe, plus la part des ventes}
\explique{Société B : $g(x)=\res{0{,}2x+800}$.}{$20\,\%=0{,}2$}
\cas{b)}
\begin{calculs}
0{,}1x+1\,200 &= 0{,}2x+800 &&\justif{même salaire dans les deux sociétés}\\
1\,200 &= 0{,}1x+800 &&\justif{on retire $0{,}1x$}\\
400 &= 0{,}1x &&\justif{on retire $800$}\\
x &= 4\,000 &&\justif{$\dfrac{400}{0{,}1}$ : diviser par $0{,}1$, c'est multiplier par $10$}
\end{calculs}
\explique{Les deux salaires sont égaux pour $\res{4\,000}$~€ de ventes.}{contrôle : $400+1\,200$
et $800+800$ font $1\,600$}
\cas{c)}
\begin{calculs}
f(5\,000) &= 0{,}1\times 5\,000+1\,200 &&\\
f(5\,000) &= 1\,700 &&\justif{$500+1\,200$}\\
g(5\,000) &= 0{,}2\times 5\,000+800 &&\\
g(5\,000) &= 1\,800 &&\justif{$1\,000+800$}
\end{calculs}
\explique{Avec $5\,000$~€ de ventes, elle gagne $100$~€ de plus dans la \res{société B} :
c'est celle-là qu'il faut choisir.}{}
\rem{Au-delà de $4\,000$~€ de ventes, B paie mieux, car elle reverse une plus grande part
des ventes ; en dessous, A paie mieux grâce à son fixe plus élevé.}
\end{exo}

\begin{exo}[Deux formules de piscine]
\cas{a)}
\explique{Formule A : $f(x)=\res{5x}$.}{pas de fixe : fonction linéaire}
\explique{Formule B : $g(x)=\res{2x+24}$.}{$24$~€ fixes, plus $2$~€ par entrée}
\cas{b) et d)}
\cotefig[0.5]{%
\explique{b) Les deux droites se coupent au point $\res{(8\,;40)}$.}{lecture en
pointillés}
\explique{d) Avant $8$ entrées, la droite de A est en dessous de celle de B : A est la
moins chère.}{la droite la plus basse donne le prix le plus bas}
\explique{Après $8$ entrées, c'est la droite de B qui est en dessous.}{}}{%
\begin{tikzpicture}[x=0.6cm,y=0.06cm]
  \draw[grille,xstep=1,ystep=10] (0,0) grid (12,60);
  \draw[->,>=Stealth,black!70] (0,0) -- (12.7,0) node[below,font=\scriptsize] {$x$};
  \draw[->,>=Stealth,black!70] (0,0) -- (0,66) node[left,font=\scriptsize] {€};
  \foreach \i in {1,...,12} {\node[below,font=\scriptsize,inner sep=1.5pt] at (\i,0) {$\i$};}
  \foreach \v in {10,20,...,60} {\node[left,font=\scriptsize,inner sep=1.5pt] at (0,\v) {$\v$};}
  \node[below left,font=\scriptsize,inner sep=1.5pt] at (0,0) {$0$};
  \draw[coulCourbe,thick] (0,0) -- (12,60);
  \draw[coulCourbeB,thick] (0,24) -- (12,48);
  \draw[lecture] (0,40) -- (8,40) -- (8,0);
  \path (8,40) node[croix]{};
  \node[coulCourbe,left,font=\small] at (10.6,54) {A};
  \node[coulCourbeB,below,font=\small] at (11.4,45) {B};
\end{tikzpicture}}
\cas{c)}
\begin{calculs}
5x &= 2x+24 &&\justif{même prix : $f(x)=g(x)$}\\
3x &= 24 &&\justif{on retire $2x$}\\
x &= 8 &&\justif{on divise par $3$}
\end{calculs}
\explique{$f(8)=5\times 8=40$ : on retrouve le point $(8\,;40)$.}{contrôle :
$g(8)=16+24=40$}
\cas{d) (conclusion)}
\explique{Pour moins de $8$ entrées, la \res{formule A} est la moins chère ; pour plus de
$8$ entrées, la \res{formule B} ; pour $8$ entrées, les deux coûtent $40$~€.}{contrôle :
$f(4)=20$ et $g(4)=32$ ; $f(12)=60$ et $g(12)=48$}
\end{exo}

\partiesansnum{Exercices type épreuve anticipée}

\begin{exo}[Automatismes (QCM)]
\cas{1.}
\explique{$10\,\%$ de $70$ font $7$, donc $30\,\%$ en font $3\times 7=21$ : réponse
\res{a)}.}{$30\,\%$, c'est $3$ fois $10\,\%$}
\cas{2.}
\explique{$100\times 1{,}2=120$, puis $120\times 0{,}8=96$ : réponse \res{b)}.}{la baisse porte
sur $120$~€ ; ou $C_{M}=1{,}2\times 0{,}8=0{,}96$}
\cas{3.}
\explique{$\dfrac{1}{8}=0{,}125$ et $0{,}125<0{,}2$ : réponse \res{c)}.}{$\dfrac{1}{4}$ et
$25\,\%$ valent $0{,}25$}
\cas{4.}
\explique{$0{,}8\times 1{,}25=1$ : réponse \res{b)}.}{à connaître : $-20\,\%$ est compensé
par $+25\,\%$}
\cas{5.}
\begin{calculs}
a &= \dfrac{7-1}{4-2} &&\justif{$\dfrac{y_{B}-y_{A}}{x_{B}-x_{A}}$}\\
a &= 3 &&\justif{$\dfrac{6}{2}$}
\end{calculs}
\explique{Réponse \res{a)}.}{}
\cas{6.}
\explique{$a=3>0$ : $f$ est croissante, réponse \res{a)}.}{le signe de $b$ ne compte pas}
\cas{7.}
\begin{calculs}
2x+3 &= 13 &&\justif{on écrit $f(x)=13$}\\
2x &= 10 &&\\
x &= 5 &&
\end{calculs}
\explique{Réponse \res{b)}.}{piège : $29$ est l'image de $13$, pas son antécédent}
\cas{8.}
\begin{calculs}
y &= -2\times 3+4 &&\justif{on remplace $x$ par $3$}\\
y &= -2 &&\justif{$-6+4$}
\end{calculs}
\explique{Réponse \res{a)}.}{}
\cas{9.}
\explique{$1{,}6\times 25=40$ : réponse \res{a)}.}{$25+0{,}6\times 25$, soit $25+15$}
\end{exo}

\begin{exo}[Les inscrits d'un club de judo]
\cas{1.}
\begin{calculs}
t &= \dfrac{171-180}{180} &&\justif{départ : 2020 ; arrivée : 2021}\\
t &= \dfrac{-9}{180} &&\\
t &= -0{,}05 &&\justif{$\dfrac{9}{180}=\dfrac{1}{20}=\dfrac{5}{100}$}
\end{calculs}
\explique{Le nombre d'inscrits a \res{baissé de $5\,\%$} entre 2020 et 2021.}{$t<0$ :
baisse}
\cas{2.}
\explique{Les inscrits sont en ligne $3$ : pour 2021, l'arrivée est en C3 et le départ
en B3.}{attention : la ligne $2$ contient les rangs}
\explique{En C4 : \res{\texttt{=(C3-B3)/B3}}.}{l'arrivée moins le départ, sur le départ}
\cas{3.}
\explique{$a=-8<0$ : $f$ est \res{décroissante}.}{propriété 4}
\explique{Selon ce modèle, le club perd \res{$8$ inscrits chaque année}.}{quand $x$
augmente de $1$, $f(x)$ diminue de $8$}
\cas{4.}
\begin{calculs}
f(8) &= -8\times 8+180 &&\justif{2028 : $x=2028-2020=8$}\\
f(8) &= 116 &&\justif{$-64+180$}
\end{calculs}
\explique{Selon ce modèle, le club compterait $\res{116}$ inscrits en 2028.}{}
\cas{5.}
\explique{On cherche d'abord quand le club a exactement $100$ inscrits.}{}
\begin{calculs}
-8x+180 &= 100 &&\justif{on écrit $f(x)=100$}\\
-8x &= -80 &&\justif{on retire $180$}\\
x &= 10 &&\justif{on divise par $-8$}
\end{calculs}
\explique{En 2030 ($x=10$), il y aurait $100$ inscrits ; $f$ est décroissante, donc
ensuite il y en a moins.}{$100$ n'est pas « moins de $100$ »}
\explique{Pour $x=11$ : $f(11)=-88+180=92$. Il y aura moins de $100$ inscrits à
partir de \res{2031}.}{$2020+11$}
\end{exo}

\begin{exo}[Les soldes]
\cas{1.}
\begin{calculs}
\text{casque} &= 80\times 0{,}75 &&\justif{$-25\,\%$ : $C_{M}=0{,}75$}\\
\text{casque} &= 60 &&\justif{les trois quarts de $80$}\\
\text{enceinte} &= 120\times 0{,}9 &&\justif{$-10\,\%$ : $C_{M}=0{,}9$}\\
\text{enceinte} &= 108 &&\justif{$120-12$}
\end{calculs}
\explique{Le casque soldé coûte $\res{60}$~€ et l'enceinte $\res{108}$~€.}{}
\cas{2.}
\explique{On paie $60+108=\res{168}$~€, au lieu de $80+120=200$~€.}{prix total avant et
après}
\begin{calculs}
t &= \dfrac{168-200}{200} &&\justif{taux d'évolution du total}\\
t &= \dfrac{-32}{200} &&\\
t &= -0{,}16 &&\justif{$\dfrac{32}{200}=\dfrac{16}{100}$}
\end{calculs}
\explique{La remise globale est de $\res{16\,\%}$.}{}
\rem{Ce n'est ni $25+10=35\,\%$, ni la moyenne $17{,}5\,\%$ : la remise de $10\,\%$ porte
sur l'article le plus cher, qui compte davantage.}
\cas{3.}
\begin{calculs}
V &= 108\times 1{,}1 &&\justif{$+10\,\%$ : $C_{M}=1{,}1$}\\
V &= 118{,}8 &&\justif{$108+10{,}8$}
\end{calculs}
\explique{$118{,}80$~€ n'est pas $120$~€ : \res{non}, le prix ne revient pas à
$120$~€.}{$0{,}9\times 1{,}1=0{,}99$, et non $1$}
\end{exo}

\end{document}
